Publish--subscribe middleware hides how failures cascade, leaving architects without pre-deployment evidence on which components are systemically critical. Software-as-a-Graph (SaG) models event-driven architectures as typed multigraphs and derives from deployment manifests an explicit dependency graph. We ask what learning-based ranking on these graphs achieves: graph neural networks on the raw multigraph and on the dependency graph, hybrids that correct a training-free baseline, and learned surrogates, evaluated on twelve synthetic architectures under leave-one-scenario-out and zero-shot on five hand-authored models of open-source systems, against three simulators. The dependency graph is what lets graph learning work: only there do messages reach the ranked components, and graph attention networks gain $+0.08$ to $+0.11$ over the raw multigraph, reaching Spearman $\rho = 0.748$. That is level with, but not equivalent to, the number of direct dependents ($0.764$), a reference that restates the reachability simulator's first propagation wave, and it rests on the degree features the networks are given. The registered primary contrast, a heterogeneous graph transformer against the training-free baseline, was null ($\Delta\rho = +0.069$, $p = 0.27$); hybrids beat the baseline ($+0.103$, $+0.130$, Holm-significant) but not their own base learners. On a queue-flow simulator whose labels cost $12.7$ CPU-hours, a gradient-boosted surrogate over dependency-graph features ranked held-out architectures above every closed-form approximation ($0.799$ against $0.706$); a graph attention surrogate did not ($0.598$). Zero-shot, the best learners reached $0.81$, but at most $0.34$ among components with nonzero impact, and their feature extraction costs $4.5$--$72\times$ one reachability simulation. Learning pays where the oracle is expensive.
